\section{Conclusion}

We introduced $(c,\epsilon)$-stationary point, a relaxed definition of Goldstein stationary point, as a new notion of convergence criterion in non-smooth non-convex stochastic optimization. Furthermore, we proposed Exponentiated O2NC, a modified online-to-non-convex framework, by setting exponential random variable as scaling factor and adopting exponentiated and regularized loss. When applied with unconstrained online gradient descent, this framework produces an algorithm that recovers standard SGDM with random scaling and finds $(c,\epsilon)$-stationary point within $O(c^{1/2}\epsilon^{-7/2})$ iterations. Notably, the algorithm automatically achieves the optimal rate of $O(\epsilon^{-4})$ for smooth objectives and $O(\epsilon^{-7/2})$ for second-order smooth objectives.


% Our results raise several interesting questions for further investigation. First, setting $c=\epsilon/\delta^2$ translates the $O(c^{1/2}\epsilon^{-7/2})$ convergence rate of our algorithm to $O(\delta^{-1}\epsilon^{-3})$, the optimal rate for finding $(\delta,\epsilon)$-Goldstein stationary point. However, whether $O(c^{1/2}\epsilon^{-7/2})$ is optimal for finding $(c,\epsilon)$-stationary point remain unclear. We conjecture that this rate is indeed optimal: an improved rate would automatically reduces to a rate better than the optimal $O(\epsilon^{-4})$ for smooth objectives.

One interesting open problem is designing an adaptive algorithm with our Exponentiated O2NC framework. Since our framework, when applied with the simplest OCO algorithm online gradient descent, yields SGDM, a natural question emerges: what if we replace online gradient descent with an adaptive online learning algorithm, such as AdaGrad? Ideally, applied with AdaGrad as the OCO subroutine and with proper tuning, Exponentiated O2NC could recover Adam's update mechanism.
% $\m_{n+1} = (1-\beta_1)\m_n + \beta_1\vecg_n$, $v_{n+1} = \beta_2v_n + (1-\beta_2)\vecg_n^2$, and $\vecx_{n+1} = \vecx_n - \eta\frac{\m_t}{\sqrt{v_t}}$. 
However, the convergence analysis for this scenario is complex and demands a nuanced approach, especially considering the intricacies associated with the adaptive learning rate.
In this vein, concurrent work by \citet{ahn2024understanding} applies a similar concept of online-to-non-convex conversion and connects the Adam algorithm to a principled online learning family known as Follow-The-Regularized-Leader (FTRL).